\chapter{嵌入式可观测性与运维}
\label{chap:observability-fleet}

实验室中可连接 JTAG、示波器和协议分析仪，现场设备通常只能提供有限日志、指标和远程诊断。可观测性必须在设计阶段建立，使开发者能够从外部输出推断系统内部状态，同时控制存储、带宽、功耗、隐私和安全风险。现场运维则把这些信号转化为资产清单、告警、分批发布、故障处置和产品改进闭环。

\section{从诊断问题而不是日志语句开始}

可观测性设计应先列出需要在没有现场仪器时回答的问题：

\begin{itemize}
  \item 设备运行的硬件、软件、配置、模型和校准版本是什么；
  \item 最近一次启动、复位、升级、回滚和模式切换发生了什么；
  \item 核心功能是否 ready，若 degraded，缺失的是哪个能力；
  \item 一次失败事务经过哪些进程、核、设备和云端阶段；
  \item 资源是瞬时峰值、持续泄漏还是下游阻塞；
  \item 故障影响一个设备、某硬件批次、某版本还是整个机群；
  \item 现场能否安全收集足够证据，并在离线状态完成恢复。
\end{itemize}

每个问题都应映射到稳定信号、保留时间、采样策略和 owner。没有明确用途的高频日志会消耗寿命和带宽；没有证据来源的支持手册则只能依赖猜测。

\section{日志、指标与追踪}

三类信号解决不同问题：

\begin{description}
  \item[日志] 记录离散事件及上下文，适合解释发生了什么；
  \item[指标] 记录 counter、gauge、histogram 和水位，适合趋势、容量与告警；
  \item[追踪] 关联一次事务跨进程、跨核和跨设备的阶段，适合定位延迟与因果关系。
\end{description}

只有自由文本日志会难以聚合，只有指标会缺少故障上下文。关键事务应同时具有 request/transaction ID、阶段指标和有限事件日志。OpenTelemetry 提供 logs、metrics、traces 与资源属性的通用语义模型\cite{opentelemetry-spec}；资源受限设备可以采用更紧凑的编码，但仍可借鉴其“资源、作用域、信号和上下文”分层。

\begin{table}[H]
  \centering
  \caption{问题与首选信号}
  \begin{tabular}{p{42mm}p{34mm}p{50mm}}
    \hline
    问题 & 首选信号 & 必要补充 \\
    \hline
    某错误首次为何发生 & 结构化事件 & 配置、版本、前后状态 \\
    P99 延迟是否恶化 & histogram/计数 & 最慢样本 trace \\
    队列是否持续积压 & depth 与 oldest age & 入/出速率、drop \\
    请求卡在哪个阶段 & 分布式或跨层 trace & 单调时间与阶段状态 \\
    重启后为何回滚 & 启动/更新事件 & 槽位元数据与复位原因 \\
    哪批设备受影响 & inventory 与标签 & 版本、硬件、地区和配置 \\
    \hline
  \end{tabular}
\end{table}

\section{事件 schema 与稳定语义}

结构化事件应包含稳定 event ID、schema 版本、严重级别、单调时间、可选实时时间、启动 ID、组件、软件版本、transaction ID 和字段。消息文本可以变化，event ID 和字段语义应保持兼容。

\begin{lstlisting}[language=C,caption={设备事件的逻辑结构}]
struct device_event {
    uint32_t schema_version;
    uint32_t event_id;
    uint64_t boot_id;
    uint64_t sequence;
    uint64_t monotonic_ns;
    int64_t  realtime_ns;
    uint32_t component_id;
    uint32_t severity;
    uint64_t transaction_id;
    uint16_t payload_length;
    uint8_t  payload[EVENT_PAYLOAD_MAX];
};
\end{lstlisting}

event ID 应表示机器可判定的事件类型，而不是源代码行号。字段要定义单位、范围、敏感级别和缺失语义。未知字段由旧解析器忽略，未知 event ID 仍应能够保留原始记录。

严重级别表示处置紧迫性，不等同于业务结果。一次重试成功的超时可以是 warning；导致安全状态的传感器失效可能是 critical。相同根因不应在每一层重复产生独立告警，可由底层记录技术事件，上层记录业务状态转换，并用 transaction ID 关联。

RFC~5424 给出了 syslog 消息格式、severity 和 structured data 的标准化思路\cite{rfc5424}，但嵌入式项目仍需为本产品定义 event catalog 和字段治理。

\section{时间、顺序与可信度}

现场证据常跨多个时钟域。每条记录应明确：

\begin{itemize}
  \item 单调时间用于持续时间和本次启动内排序；
  \item 实时时间用于跨启动和外部事件关联，但可能未同步或回退；
  \item boot ID/启动计数把相同单调时间范围分开；
  \item sequence 检测 ring 覆盖、丢包和乱序；
  \item 时间同步状态、offset 与不确定度；
  \item 服务器接收时间只能作为外部边界，不能替代设备事件时刻。
\end{itemize}

首次联网前的日志不应伪装成可信日历时间。可以记录“实时时间未知”，并在后续上传时携带单调时间、启动 ID 和首次可信时间映射。跨设备时序分析见第~\ref{chap:embedded-network-time}~章。

\section{高频路径与观测开销}

ISR、fault handler 和硬实时路径不应执行格式化、文件 I/O、锁等待或网络发送。可以写入固定长度的 per-CPU/ring buffer，再由后台任务编码和持久化。

观测系统本身需要预算：

\begin{itemize}
  \item 每类事件最大速率和字节数；
  \item ring 容量、覆盖策略与保留窗口；
  \item CPU 格式化成本、锁竞争和 cache 影响；
  \item Flash 写入量、写放大和每日寿命预算；
  \item 无线唤醒、上传流量和每字节能量；
  \item debug 模式的最大持续时间和自动恢复。
\end{itemize}

高频事件可使用采样、聚合、token bucket 或“首次 + 计数 + 周期摘要”。限流本身也要产生丢弃计数，不能静默丢失。安全、崩溃和更新结果可获得更高优先级，但仍需有界。

\section{指标设计与基数控制}

指标应直接对应产品 SLO、资源预算和故障假设，例如：

\begin{itemize}
  \item 任务/线程最大延迟、队列水位、oldest age 和丢弃计数；
  \item 内存、fd、DMA buffer、CMA 和存储剩余；
  \item ECC corrected/failed、I/O error 和文件系统只读切换；
  \item 网络重连、丢包、RTT 和时间同步 offset；
  \item 温度、降频、功耗状态、唤醒原因和电池估计；
  \item watchdog、崩溃、回滚、启动失败和安全状态次数。
\end{itemize}

仅报告当前值会遗漏峰值，关键资源应保存 high-water mark 和 histogram。计数器需要定义重启、回绕、持久化和清零权限。

标签会增加时序数量。把 request ID、完整 URL、用户 ID 或任意错误文本作为指标标签会形成高基数，耗尽设备和后端资源。高基数字段应进入有限日志或 trace；指标标签只保留稳定、有限的版本、硬件、地区、结果分类和模式。

histogram bucket 应围绕产品门槛设计。例如要求 P99 小于 100\,ms 时，需要在 100\,ms 附近具有足够分辨率。只有平均值和最大值无法判断尾部变化。

\section{跨层事务追踪}

控制命令、媒体帧、OTA 和生产操作可以使用 transaction ID 关联用户空间、内核、远端核和云端事件。W3C Trace Context 定义了 \texttt{traceparent} 与 \texttt{tracestate} 的跨系统传播格式\cite{w3c-trace-context}；本地二进制协议可以使用更紧凑字段，但应明确 trace ID、span/阶段 ID、采样标志和安全边界。

跨层追踪应围绕状态边界：

\begin{itemize}
  \item 请求进入、排队、开始服务和返回；
  \item 内核 ioctl/提交、DMA、IRQ 和完成；
  \item 远端核接收、执行与应答；
  \item 设备真实采样或执行器生效；
  \item 网络发送、服务端接受与业务确认。
\end{itemize}

不要给每个函数打印 enter/exit。高频路径使用采样、硬件 trace 或只保留超阈值样本。采样决定应尽可能在一次 trace 内一致，否则只保留末端 span 会丢失因果。

外部传入 trace ID 属于不可信字段，必须校验格式并限制其作为索引和日志内容的资源影响。内部安全事件不应仅凭调用方 trace 标识建立授权。

\section{日志容量、持久性与丢失语义}

日志必须有速率、容量、优先级、轮转和保留策略。建议分层：

\begin{description}
  \item[RAM ring] 高频调试与最近上下文，重启后可丢失或由保留内存保存；
  \item[持久运行日志] 有界轮转，用于跨重启故障分析；
  \item[审计日志] 记录高价值安全和管理动作，访问与完整性更严格；
  \item[上传队列] 保存尚未送达的摘要、指标和关键事件；
  \item[崩溃区] 写入路径极小、抗二次故障、由后续启动提取。
\end{description}

“写入日志 API 成功”可能只表示进入内存队列，不代表已持久。对于审计和关键更新结果，要定义所需持久边界；对于普通 debug，过度 fsync 会损害性能和 Flash 寿命。

存储满时必须保留核心控制和恢复能力。日志系统应预留空间或使用独立配额，优先淘汰低价值旧记录，并记录覆盖计数。不可让日志写失败阻塞实时线程或触发无限递归日志。

\section{崩溃证据与符号化}

MCU HardFault 应保存寄存器、堆栈、任务、复位原因、固件版本和有限内存窗口；Linux 保存 pstore/ramoops、core dump、内核 oops 和服务状态。Linux pstore 文档描述了跨重启保存平台崩溃记录的接口\cite{linux-pstore}，具体容量和后端取决于平台。

崩溃区写入必须有界且抗二次故障。若 fault handler 再次异常，应保留首次故障最小头部，而不是反复覆盖。建议包含：

\begin{itemize}
  \item magic、schema、长度、CRC、generation 与完整标志；
  \item build ID、镜像摘要、硬件修订和启动 ID；
  \item 异常类型、PC/SP、寄存器、线程和 CPU；
  \item 最近固定数量的关键事件和资源水位；
  \item 是否发生栈截断、读取失败或二次 fault。
\end{itemize}

所有地址解析都依赖完全匹配的 ELF、内核模块、映射和符号。符号服务应按 build ID 保存未剥离产物、调试文件、map、配置和源码提交。systemd-coredump 等机制可以管理用户空间 core，但其存储、压缩、权限和敏感数据策略必须按产品配置\cite{systemd-coredump}。

崩溃上传成功后不应立即删除唯一原始证据。采用“已提取、已确认、可回收”的状态，并限制重复上传。

\section{现场诊断包}

诊断包可以包含：

\begin{itemize}
  \item 硬件、软件、配置、容器、模型和校准 manifest；
  \item 启动、复位、更新、槽位和安全状态摘要；
  \item 有界日志、指标 snapshot、最慢 trace 和崩溃记录；
  \item 进程/线程、cgroup、fd、mount、网络和时间同步状态；
  \item 存储健康、文件系统、温度、电源和设备错误；
  \item 采集工具版本、开始/结束时间、失败项和每个文件摘要。
\end{itemize}

生成过程必须设置大小、CPU、I/O 和总时间上限，并允许部分成功。每个 collector 独立超时；一个挂死驱动不能阻止收集其他证据。

\begin{lstlisting}[language=C,caption={诊断包条目的逻辑清单}]
struct diagnostic_item {
    uint32_t collector_id;
    uint32_t schema_version;
    uint64_t collected_monotonic_ns;
    uint64_t size_bytes;
    uint8_t  sha256[32];
    int32_t  status;
    uint32_t sensitivity;
};
\end{lstlisting}

敏感密钥、token、用户内容和完整内存不应默认进入诊断包。字段应分级，导出前执行脱敏、访问授权和加密封装。对无法安全自动脱敏的 core dump，可只在本地授权维修或专用安全通道中收集。

\section{机群 inventory 与状态模型}

远程设备管理至少需要可靠 inventory：

\begin{itemize}
  \item 设备身份、产品、硬件修订、关键器件批次和地区；
  \item Bootloader、固件、内核、rootfs、应用、模型和配置版本；
  \item 当前槽、最后已知良好版本、更新能力和证书状态；
  \item 在线/离线、最近联系、健康、降级和支持期；
  \item 生产、维修、所有权和退役生命周期状态。
\end{itemize}

普通应用上报的版本字符串可能错误或被篡改。高价值决策宜与签名 manifest、受保护启动记录或证明结果核对。

控制面应区分期望状态（desired）与设备确认状态（reported）。云端写入 desired 只表示意图被接受，不表示设备已经执行。每次状态变更携带 generation、适用条件和结果，防止离线设备重连后执行陈旧命令。

NIST IoT 核心能力基线强调设备识别、配置、数据保护、接口限制、软件更新和状态感知等能力\cite{nist-iot-core-baseline}；这些能力只有被 inventory 与运维流程持续观测，才能在现场形成闭环。

\section{远程命令与作业状态机}

任意 shell 不应作为常规运维接口。诊断、重启、抓包、日志提升和证书轮换应建模为受限命令，每条包含：

\begin{itemize}
  \item command ID、类型、目标设备和参数 schema；
  \item 发起者、审批、权限和审计上下文；
  \item 创建时间、过期时间和适用设备条件；
  \item 幂等键、最大执行次数和取消规则；
  \item queued、accepted、running、succeeded、failed、expired 等状态；
  \item 结果摘要、证据和稳定错误分类。
\end{itemize}

设备应在本地再次检查模式、版本、电量、物理安全条件和资源，不应因服务端已授权就执行明显危险操作。网络重试必须返回同一命令的既有结果，避免重复重启、擦除或执行器动作。

长任务需要心跳和进度，但进度不能替代最终业务确认。设备断线后服务端应保持“结果未知”，而不是自动标记失败并重新创建无幂等保障的新任务。

\section{分批发布、错误预算与回滚}

OTA 应按实验设备、内部设备、小比例 canary 和逐步扩大批次发布。切片不仅按随机百分比，还应覆盖硬件修订、存储料号、地区、网络、使用强度和关键客户场景。

每一阶段设置自动暂停门槛：

\begin{itemize}
  \item 下载、验证、安装、试启动和健康确认成功率；
  \item 启动时间、崩溃、watchdog、回滚和恢复率；
  \item 核心功能错误、P99 延迟、内存、温度和功耗；
  \item 安全事件、证书、数据迁移和存储健康；
  \item 与匹配基线批次相比的统计差异。
\end{itemize}

错误预算应同时包含分母和窗口。例如“候选版本 24 小时内启动失败率高于 0.2\% 且至少 20 台设备失败则暂停”，比单个绝对失败数更能适应批次规模。稀有严重故障仍可设置“一票暂停”。

设备回滚不等于云端配置、模型和数据自动回滚。发布系统需要记录固件、服务端能力、配置和 schema 的兼容关系，并验证长期离线设备跨多个版本升级。

\section{离线、弱网与 store-and-forward}

现场设备可能长期离线。遥测队列应有容量、优先级和淘汰策略，安全事件、升级结果和崩溃摘要不能被大量普通指标挤出。

每类记录定义：

\begin{itemize}
  \item 最大本地保留时间和字节预算；
  \item 是否可聚合、降采样或只保留最新值；
  \item 丢弃顺序和丢失计数；
  \item 重连后的上传速率、批大小和压缩；
  \item 服务端去重键与乱序处理；
  \item 电量低、计费链路和业务拥塞时的行为。
\end{itemize}

重连后不能一次上传全部积压而占满业务链路或持续唤醒无线模块。采用带随机抖动的速率控制，并让核心控制和更新流量具有独立配额。

设备本地必须能够完成核心诊断和恢复，不应把云端在线作为进入安全状态、回退或读取复位原因的必要条件。

\section{告警与事件关联}

告警应对应可采取动作的状态，而不是每条错误日志。一个稳定告警至少包含：

\begin{itemize}
  \item 触发条件、窗口、迟滞和恢复条件；
  \item 影响范围、严重度、owner 和升级路径；
  \item 关联 runbook、诊断查询和自动收集动作；
  \item 去重、抑制和父子依赖；
  \item 何时属于设备故障、基础设施故障或未知。
\end{itemize}

例如大量设备同时报告 DNS 失败，应优先关联地区/运营商/云端变更，而不是为每台设备创建独立工单。相反，单批次存储 ECC 增长需要按硬件 lot 切片，即使总体机群比例很低。

告警恢复不表示根因已关闭。短暂重启可能让指标恢复，但崩溃证据仍需要归档和问题追踪。

\section{事故响应与证据闭环}

现场事故处理建议采用：

\begin{enumerate}
  \item 确认用户影响和安全边界，必要时停止发布或限制功能；
  \item 冻结受影响版本、设备切片和原始遥测；
  \item 建立统一事件时间线与假设列表；
  \item 通过诊断包、符号化和实验室复现缩小故障层次；
  \item 发布有界缓解、修复和恢复方案；
  \item 验证机群覆盖、残余风险和是否存在滞留离线设备；
  \item 完成根因、逃逸原因、监测缺口与系统性改进。
\end{enumerate}

根因分析应区分触发条件、代码/硬件缺陷、逃逸原因和放大因素。修复一个空指针还不够；还应补充输入契约、回归测试、崩溃分类和发布门，防止同类问题再次逃逸。

现场事件应反馈到需求、风险、HIL 场景和设计预算。可观测性若只帮助“解释过去”，却不改变测试与架构，就没有完成工程闭环。

\section{隐私与数据治理}

遥测应遵循最小收集、目的限定和保留期限。设备唯一身份、位置、无线邻居、音视频、用户操作和故障内存可能属于敏感数据。

对每个字段维护数据字典：

\begin{itemize}
  \item 采集目的、合法依据和必要性；
  \item 数据类型、精度、采样率和去标识方法；
  \item 设备、本地网关和云端各自保留时间；
  \item 访问角色、导出、跨区域和第三方共享；
  \item 用户删除、设备转让和退役时的处理；
  \item 是否可能与其他字段组合重新识别个人。
\end{itemize}

脱敏应在尽量靠近数据源的位置执行。日志进入上传队列后再做正则替换，可能已经把秘密写入 Flash。结构化 schema 允许按字段执行 redact、hash、truncate 或禁止采集，比自由文本更可靠。

\section{运维安全与维修模式}

远程诊断、日志等级、抓包、core dump 和固件下载都会扩大攻击面。操作需要最小权限、短期授权、审计、速率限制和设备端策略。

“break-glass” 紧急权限应：

\begin{itemize}
  \item 只针对特定设备、操作和有限时间；
  \item 需要强认证、审批与不可变审计；
  \item 在设备端显示或记录紧急模式；
  \item 自动到期并撤销临时凭据；
  \item 完成后验证调试口、账户和服务恢复量产策略。
\end{itemize}

现场维修模式宜由物理动作、设备生命周期状态和签名授权共同开启。服务端凭据泄漏时，设备端仍应拒绝不满足本地安全条件的擦除、密钥导出和执行器命令。

\section{数字化验收矩阵}

\begin{table}[H]
  \centering
  \caption{可观测性与机群运维验收矩阵示例}
  \begin{tabular}{p{28mm}p{41mm}p{39mm}p{32mm}}
    \hline
    场景 & 条件与注入 & 关键证据 & 通过条件 \\
    \hline
    日志风暴 & 单一错误高频触发 & CPU、写入、drop、寿命 & 有界聚合且核心业务不退化 \\
    崩溃循环 & 多阶段重复崩溃 & 首次 core、pstore、build ID & 证据不被覆盖且可符号化 \\
    长期离线 & 队列满后重连 & 优先级、丢失计数、带宽 & 关键事件保留且限速收敛 \\
    远程命令 & 重放、过期、断线与取消 & 审计、状态机和设备状态 & 无重复危险动作 \\
    分批发布 & 指定硬件切片注入故障 & 基线比较、暂停和回滚 & 门槛触发后停止扩大 \\
    隐私检查 & 日志、包、core 与删除 & 字段清单、扫描和审计 & 无未批准敏感数据残留 \\
    \hline
  \end{tabular}
\end{table}

验收应在真实存储容量、弱网、电量、温度和高业务负载下执行。测试构建中无限磁盘和高速网络得到的结果不能证明量产策略有效。

\section{可观测性与运维检查清单}

\begin{itemize}
  \item 每个现场诊断问题是否具有明确日志、指标、trace 或 inventory 来源？
  \item event ID、字段、单位、敏感级别和兼容策略是否稳定？
  \item 单调时间、实时时间、boot ID、sequence 和不确定度是否区分？
  \item 高频路径是否使用有界 ring，并量化 CPU、存储、带宽和功耗开销？
  \item 指标是否保存水位与 histogram，并限制标签基数？
  \item 跨层事务是否能关联应用、内核、远端核、设备和云端阶段？
  \item 日志满、上传队列满和持续离线时是否有明确丢失语义？
  \item 崩溃信息是否能按 build ID 找到完全匹配的符号和配置？
  \item 诊断包是否有界、可部分成功并执行敏感数据过滤？
  \item 远程命令是否具有 ID、过期、授权、审计和幂等状态机？
  \item OTA 是否按风险切片发布并由健康指标自动暂停？
  \item 告警是否可操作、可去重，并区分设备与基础设施故障？
  \item 事故结论是否反馈到需求、测试、HIL、预算和发布门？
  \item 维修与紧急权限是否设备绑定、自动到期并可审计？
\end{itemize}
